\documentclass{article}

\textwidth=6.5in
\textheight=8.5in
\voffset=-54pt
\oddsidemargin=0in
\evensidemargin=0in
\setlength{\parskip}{8pt}
\setlength{\parindent}{0pt}

\usepackage{amsmath, amssymb}
\usepackage{ifthen}

\usepackage[inline]{enumitem}
\setlist{topsep=0pt, parsep=8pt}

\usepackage[rgb,table]{xcolor}\convertcolorsUfalse
\usepackage{graphicx}

% change hyperref options
\PassOptionsToPackage{hyphens}{url}
\usepackage[bookmarks,colorlinks,linkcolor=black,citecolor=blue,pdfstartview=FitH,urlcolor=blue]{hyperref}
\hypersetup{pdfpagemode=UseNone}
\usepackage{bookmark}

\usepackage{graphicx}
\usepackage{multicol}
\usepackage{framed}
\usepackage{array}

\usepackage{icomma}

\usepackage{soul}

\usepackage{tikz}
\usetikzlibrary{
  matrix,%
  % enables the snake paths
  decorations.pathmorphing,%
  % enables bigger arrows
  decorations.markings,%
  decorations.pathreplacing,%
  decorations.text,%
  calc,%
  patterns,%
  shapes.geometric,%
  arrows,
  decorations.shapes,
  positioning,
  plotmarks,
  intersections
}
\tikzset{>=latex} % sets default arrow type in tikz
\tikzset{font=\normalsize}
\tikzset{mark size=1.5pt, mark options=thin}
\tikzset{pin distance=4pt,
  every pin edge/.style={<-, thin, shorten <= -2pt}}

\interfootnotelinepenalty=10000
\renewcommand{\thefootnote}{(\arabic{footnote})}

\usepackage{multirow, bigdelim}

\usepackage[notref,notcite]{showkeys}

% change to \usepackage[sols]{handout} to see solutions
\usepackage[sols,notes,workshop,grading]{handout}

\newcommand\iscurrentenumi[1]{%
  \ifthenelse{\equal{#1}{\theenumi}}{}{#1}}
\setenumerate[1]{ref=\#\arabic*}
\setenumerate[2]{ref=\protect\iscurrentenumi{\theenumi}(\alph*)}
\newcommand\enumiiprefix[2]{%
  \ifthenelse{{\equal{#1}{\theenumi}}\AND{\equal{#2}{(\alph{enumii})}}}{}{}%
  \ifthenelse{{\equal{#1}{\theenumi}}\AND{\NOT{\equal{#2}{(\alph{enumii})}}}}{#2}{}%
  \ifthenelse{\equal{#1}{\theenumi}}{}{#1#2}%
}
\setenumerate[3]{ref=\protect\enumiiprefix{\theenumi}{(\alph{enumii})}\roman*}

\definecolor{purple}{rgb}{0.5,0,1.0}

\usepackage{fancyhdr}
\lhead{\textcolor{white}{This content is protected and may not be shared, uploaded, or distributed.}}
\rhead{}
\rfoot{\vspace*{\fill}\footnotesize\textcopyright{} \emph{Harvard Math Ma}}
\renewcommand{\headrulewidth}{0pt}
\pagestyle{fancy}
\begin{document}

\begin{center}
  \Large\textsc{Exam 1 Learning Objectives}
\end{center}

\begin{enumerate}
\item \textbf{Introduction to Functions} %% 1
  \begin{itemize}

  \item Explain what a \emph{function} is, and distinguish between functions and non-functions.

  \item Identify the \emph{independent variable}, \emph{dependent variable}, \emph{domain}, and \emph{range} of a function.

  \item Explain how slope graphically represents a rate of change.

  \item Reason about whether a function is \emph{increasing} or \emph{decreasing}, as well as whether it's \emph{concave up}, \emph{concave down}, or \emph{linear}.

  \end{itemize}

\item \textbf{A Library of Basic Functions and Properties} %% 2
  \begin{itemize}
  \item Sketch and recognize graphs of $y = mx + b$, $y = x^2$, $y = x^3$, $y = \frac{1}{x}$, $y = \frac{1}{x^2}$, $y = |x|$, and $y = \sqrt{x}$.

  \item Understand the following characteristics of a function graphically: positive / negative, increasing / decreasing, intercepts, asymptotes, continuity, concavity, even / odd.

  \item Algebraically determine whether a function is even, odd, or neither.

  \item Explain what it means to say a function is increasing / decreasing \emph{without bound}, as opposed to having an \emph{upper bound} or \emph{lower bound}.

  \item Recognize when the graph of a function given algebraically has a hole and construct algebraic examples of functions whose graphs have holes.

  \end{itemize}

\item \textbf{More on Functions, Average Rates of Change} %% 3
  \begin{itemize}

  \item Write and interpret expressions in function notation.

  \item Compute an average rate of change, interpret it in the context of a situation, and represent it graphically.

  \item Explain the difference between average speed, average velocity, instantaneous speed, and instantaneous velocity.

  \item Explain how concavity can be characterized in terms of secant line segments.

  \end{itemize}

\item \textbf{Transforming Functions} %% 4
  \begin{itemize}
  \item Sketch $y = a f(bx) + c$, $y = a f(x + b) + c$, and $y = |f(x)|$
    from a graph of $y = f(x)$.

  \end{itemize}

\item \textbf{Linear and Piecewise Linear Functions} %% 5
  \begin{itemize} 
  \item Write the equation of a line given 2 points or 1 point and the
    slope, and decide whether slope-intercept form ($y = mx + b$) or
    point-slope form is more useful. 

  \item Use linear and piecewise linear equations to model various
    phenomena.

  \item Interpret the slope in linear and piecewise linear functions, 
    and graph the ``slope function'' $f'(x)$ of a piecewise linear
    function $f(x)$.  

  \end{itemize}

\item \textbf{Function Composition, Modeling with Functions} %% 6
  \begin{itemize}
  \item Explain how composition of functions models situations where there is a chain of dependency among variables.

  \item Compute compositions of functions.

  \item Model situations systematically: identify your goal, define variables and express your goal in terms of those variables, and then look for relationships among your variables.

  \item Communicate your modeling process clearly. 
  \end{itemize}
\end{enumerate}

\begin{enumerate}[start=8, topsep=0pt]
\item \textbf{Approximating with Secant Lines} %% 8
  \begin{itemize}
  \item Decide where the graph of a function is or isn't locally linear.

  \item Use secant lines to approximate function values.

  \item Determine whether a secant line approximation is too high or too low, or if there isn't enough information to be sure.

  \end{itemize}

\item \textbf{Definition and Interpretation of the Derivative} %% 9
  \begin{itemize}
  \item State the definition of the derivative, explain where it comes from, and use it to calculate derivatives.

  \item Interpret the derivative in context and represent it graphically.

  \item Find the tangent line to the graph of $f(x)$ at $x = a$. 

  \item Use tangent lines to approximate function values, and determine whether such an approximation is too high or too low, or if there isn't enough information to be sure.
  \end{itemize}

\item \textbf{The Derivative Function} %% 10
  \begin{itemize}
  \item Use the limit definition to compute the derivative function
    $f'$.

  \item Determine whether a function is \emph{differentiable}, both
    graphically and algebraically, and give examples of functions that
    are differentiable and non-differentiable.

  \item Sketch a rough graph of $f'$ when you have the graph of $f$. 

  \end{itemize}

\item \textbf{Functions and their Derivatives} %% 11
  \begin{itemize}

  \item Explain how concavity can be characterized in terms of tangent lines.

  \item Explain what $f'$ and $f''$ tell us about the graph of $f$. (For example, if $f'$ is positive, what does that say about $f$? If $f'$ is increasing, what does that say about $f$?)

  \item Explain how displacement, velocity, and acceleration are related in terms of derivatives.

  \end{itemize}

\item \textbf{Using Derivatives to Graph Functions} %% 12
  \begin{itemize}

  \item Use the first derivative to determine where a function is
    increasing or decreasing and to find \emph{turning points}.

  \item Use the second derivative to determine where a function is
    concave up or concave down and to find \emph{inflection points}. 

  \item Put the above information together to sketch a reasonable graph
    of a function.

  \item Develop the habit of checking for consistency when putting several pieces of information together. 
  \end{itemize}

\end{enumerate}
\end{document}
